% APÊNDICES

\begin{apendicesenv}
%\partapendices % Folha divisória (comentar caso não deseje que seja exibida)

% Primeiro apêndice
\chapter{Nome do primeiro apêndice} % Edite para alterar o título deste apêndice
\label{chap:apendiceA}

Os anexos são documentos complementares e/ou comprobatórios do texto da dissertação. São itens opcionais.

\begin{adjustwidth}{4cm}{}
    \SingleSpacing
    \small
    
    Trazem informações esclarecedoras que não se incluem no texto para não prejudicar a sequência lógica da leitura. O apêndice difere do anexo por ser elaborado pelo próprio autor, enquanto o anexo é um documento de autoria de outro(s). Segundo a NBR 6029 (ABNT, 2006), tanto o apêndice quanto o anexo são identificados por letras maiúsculas sequenciais, seguidos de seus respectivos títulos. (FRANÇA; VASCONCELLOS, 2013, p. 23).
    
    \vspace{0.4cm}
\end{adjustwidth}

Os itens são identificados por letras maiúsculas sequenciais, travessão e seguidos de seus respectivos títulos. Tal título deve ser colocado no alto da página, centralizado, escrito em Arial ou Times New Roman, tamanho 12, espaçamento 1,5. Depois, inserir o conteúdo do apêndice ou anexo. Devem ser inseridos após a lista de referências bibliográficas, sendo primeiro os apêndices e, por últimos, os anexos. A paginação deve dar sequência à numeração do texto.

São considerados apêndices: Formulários e questionários aplicados ou o roteiro da entrevista; Planos de ensino e de aula, criados para a aplicação da metodologia proposta; Regulamentos e regras criados para a implantação do projeto-piloto.

Lembre-se que a diferença entre apêndice e anexo diz respeito à autoria do texto e/ou material ali colocado.

Caso o material ou texto suplementar ou complementar seja de sua autoria, então ele deverá ser colocado como um apêndice. Porém, caso a autoria seja de terceiros, então o material ou texto deverá ser colocado como anexo.

Caso seja conveniente, podem ser criados outros apêndices/anexos para o seu trabalho acadêmico. Basta recortar e colar este trecho neste mesmo documento. Lembre-se de alterar o "label".

% Segundo apêndice
\chapter{Nome do segundo apêndice}
\label{chap:apendiceB}

Conteúdo do segundo apêndice

\end{apendicesenv}
